\subsection{Lagrangian subspaces}

For a symplectic space $(V,\omega)$, let $\L(V,\omega)$ denote the set of its Lagrangian subspaces.
We denote
$$ \L(n) = \L(\bR^{2n},\omega_0) . $$
We can characterize the Lagrangian subspaces of $\bR^{2n}$ more explicitly.
Firstly a Lagrangian subspace of $\bR^{2n}$ has dimension $n$, and is thus the image of some full-rank linear transform $\bR^n\to\bR^{2n}$, or equivalently a $2n\times n$-matrix of rank $n$.
We can then determine when such an image is Lagrangian.

\bprop

    Let $X,Y\in\Mat_n$ be matrices, and define $\Lambda\leq\bR^{2n}$ by
    $$ \Lambda = \im Z,\qquad Z=\pmatrix{X\cr Y} . $$
    Then $\Lambda\in\L(n)$ iff $Z$ has rank $n$ and $X^\top Y=Y^\top X$.

    In particular the graph $\Lambda=\set{(x,Ax)}[x\in\bR^n]$ of some matrix $A\in\Mat_n$ is Lagrangian iff $A$ is symmetric.

\eprop

\bproof

    It is sufficient to consider when $\omega_0$ is trivial on $\Lambda$.
    Let $z=(Xu,Yu)$ and $z'=(Xu',Yu')$ in $\Lambda$, then
    $$ \omega_0(z,z') = \pmatrix{(Xu)^\top&(Yu)^\top}\pmatrix{0&-I\cr I&0}\pmatrix{Xu'\cr Yu'} = u^\top(X^\top Y-Y^\top X)u' . $$
    This is zero for all $u,u'$ iff $X^\top Y=Y^\top X$, as desired.

    The second claim is because $\Lambda=\im(I,A)$, which has rank $n$.
    Thus it is Lagrangian iff $A^\top=A$.
    \qed

\eproof

A matrix $Z\in\bR^{2n\times n}$ of the form
$$ Z = \pmatrix{X\cr Y},\qquad X^\top Y=Y^\top X $$
is called a {\emphcolor Lagrangian frame}.
Notice that $Z$'s columns form an orthonormal basis for $\Lambda=\im Z$ iff $U=X+iY$ is unitary.
Indeed $Z$'s columns are orthonormal iff
$$ I = Z^\top Z = X^\top X + Y^\top Y . $$
And $U$ is unitary iff
$$ I = U^\dagger U = (X^\top - iY^\top)(X+iY) = X^\top X + Y^\top Y + i(X^\top Y - Y^\top X) = X^\top X + Y^\top Y . $$
A Lagrangian frame $Z$ for which $U=X+iY$ is unitary is called a {\emphcolor unitary Lagrangian frame}.
Every Lagrangian subspace has a unitary Lagrangian frame as we can take an orthonormal basis and put it in a matrix.

Let us define the Lagrangian subspace
$$ \Lambda_\hor = \set{(x,0)\in\bR^{2n}}[x\in\bR^n] = \im\pmatrix{I\cr0} . $$

\blemm

    \benum
        \item If $\Lambda\in\L(n)$ and $\Psi\in\Sp(2n)$, then $\Psi\Lambda\in\L(n)$.
        \item For any two Lagrangian subspaces $\Lambda,\Lambda'\in\L(n)$, there is a symplectic matrix $\Psi\in\unitary(n)$ such that $\Lambda'=\Psi\Lambda$.
        \item $\L(n)$ can be seen to be naturally isomorphic to $\unitary(n)/\orth(n)$, and thus carries the natural structure of a manifold of dimension $n(n+1)/2$.
    \eenum

\elemm

\bproof

    (1) is clear.
    For (2) let $\Lambda\leq\bR^{2n}$ be given by a unitary Lagrangian frame $Z=(X,Y)$ and define
    $$ \Psi = X + iY = \pmatrix{X&-Y\cr Y&X} \in \Sp(2n)\cap\orth(n) . $$
    Then
    $$ \Psi\Lambda_\hor = \Lambda , $$
    and thus we have shown (2).

    A unitary Lagrangian frame of $\Lambda$ is uniquely determined up to choice of orthonormal basis.
    Altering the basis amounts to multiplication on the right by an orthonormal matrix, so we see that there is a bijection $\unitary(n)/\orth(n)\cong\L(n)$.
    The dimension of $\unitary(n)$ is $n^2$, and the dimension of $\orth(n)$ is $n(n-1)/2$, so $\dim\L(n)=n(n+1)/2$.
    \qed

\eproof

We can show that $\L(n)\cong\unitary(n)/\orth(n)$ has a fundamental group isomorphic to the integers.
This isomorphism is also called the {\emphcolor Maslov index}.

\bthrm

    There exists a unique function, the {\emphcolor Maslov index}, $\mu$ which assigns to each loop $\Lambda\colon\bR/\bZ\to\L(n)$ an integer and satisfies the following axions:
    \benum
        \item homotopy: two loops in $\L(n)$ are homotopic iff they have the same Maslov index.
        \item product: for any two loops $\Lambda\colon\bR/\bZ\to\L(n)$ and $\Psi\colon\bR/\bZ\to\Sp(2n)$, we have
        $$ \mu(\Psi\Lambda) = \mu(\Lambda) + 2\mu(\Psi) . $$
        \item direct sum: if $n=n'+n''$, then we can identify $\L(n')\oplus\L(n'')$ as a submanifold of $\L(n)$ and then
        $$ \mu(\Lambda'\oplus\Lambda'') = \mu(\Lambda') + \mu(\Lambda'') . $$
        \item zero: a constant loop $\Lambda(t)=\Lambda_0$ has Maslov index zero.
        \item normalization: the loop $\Lambda\colon\bR/\bZ\to\L(1)$ given by
        $$ \Lambda(t) = e^{\pi it}\bR \leq \bC = \bR^2 $$
        has Maslov index one.
    \eenum

\ethrm

\bproof

    An orthonormal matrix has determinant $\pm1$, thus the map
    $$ \rho\colon\L(n)\cong\unitary(n)/\orth(n) \to S^1,\qquad \rho([U]) = \det_\bC U^2 $$
    is well-defined.
    Here $[U]\in\unitary(n)/\orth(n)$ denotes the coset of the unitary matrix $U=X+iY$, which in turn represents the Lagrangian subspace $\Lambda=\im(X,Y)$.

    Define the Maslov index to be the induced map $\pi_1(\L(n))\to\pi_1(S^1)\isoto\bZ$.
    That is for a loop $\Lambda$,
    $$ \mu(\Lambda) = \deg(\rho\circ\Lambda) = \alpha(1) - \alpha(0) $$
    where $\alpha\colon I\to\bR$ is a lift of $\rho\circ\Lambda$, meaning
    $$ e^{\pi i\alpha(t)} = \det_\bC(U(t)),\qquad \Lambda = [U] . $$

    The Maslov index then depends only on the homotopy class of $\Lambda$.
    We show it satisfies the product axiom: first note that if $\Psi$ is unitary $\rho(\Psi\Lambda)=\rho(\Psi)^2\rho(\Lambda)$ ($\rho(\Psi)$ corresponds to the $\rho$ in \refmath[theorem]{maslovsp}).
    And so the degree of this is
    $$ \deg(\rho(\Psi\Lambda)) = \deg\rho(\Lambda) + 2\deg\rho(\Psi) = \mu(\Lambda) + 2\mu(\Psi) . $$
    If $\Psi$ is not unitary, it is homotopic to a loop of unitary matrices, and the result follows.

    We now show the converse direction in the homotopy axiom.
    Suppose that $\Lambda_0,\Lambda_1$ are two loops of Lagrangian subspaces with the same Maslov index, we must show $\Lambda_0,\Lambda_1$ are homotopic.
    By the previous lemma and the product axiom, we may assume $\Lambda_j(0)=\Lambda_j(1)=\Lambda_\hor$.
    Lift $\Lambda_j(t)$ to $U_j(t)=X_j(t)+iY_j(t)\in\unitary(n)$ such that $U_0(0)=U_1(0)=I$ (as these represent $\Lambda_\hor$).
    We may multiply $U_j$ by right multiplication with orthogonal matrices to get $U_0(1)=U_1(1)=I$ as well.
    Then $U(t)=U_1(t)U_0(t)^{-1}$ forms a loop of unitary matrices.

    Note that $\mu(\Lambda_0)=\mu(\Lambda_1)$ means $2\mu(U_0)=2\mu(U_1)$, and thus $U_0,U_1$ are homotopic.
    Thus $\Lambda_0,\Lambda_1$ are homotopic.

    The normalization axiom follows from the other axioms.
    Note that the loop $\Lambda(t)=e^{\pi it}\bR$ satisfies
    $$ \Lambda(t) \oplus \Lambda(t) = \Psi(t)\Lambda_0 $$
    where $\Lambda_0=\bR^2\leq\bC^2$ is constant and $\Psi$ is the loop
    $$ \Psi(t) = e^{2\pi it}X(t) = e^{\pi it}\pmatrix{\cos(\pi t)&-\sin(\pi t)\cr \sin(\pi t)&\cos(\pi t)} . $$
    Here $X(t)$ has constant determinant $1$ and so $\det\Psi=e^{2\pi it}$ lifts to $\alpha(t)=t$ and thus $\mu(\Psi)=1$.
    Thus we have
    $$ 2\mu(\Lambda) = \mu(\Lambda\oplus\Lambda) = \mu(\Lambda_0) + 2\mu(\Psi) = 2 \implies \mu(\Lambda) = 1 . \qed $$

\eproof

